% GENERATED FILE. Edit canonical lessons, then run npm run generate:books.
% canonical-source-hash: fnv1a64:8b3ec98e779b45a3
% canonical-lessons: AR-C213-arika, AR-C213-fursha, AR-C213-ghita, AR-C213-athath, AR-C213-ghasil

\chapter{Sofa, Brush, Lid, Furniture, Laundry}
\label{ch:ar-sofa-brush-lid-furniture-laundry}
\hlchaptermodality{213}

\begin{chapteropening}
\textbf{By the end of this chapter:} \emph{I can say a sofa, a brush, a lid, furniture and laundry.}

Complete the last lesson of chapter 213: I can say a sofa, a brush, a lid, furniture and laundry.
\end{chapteropening}

\section[\texorpdfstring{\ar{أريكة} (arīka)}{أريكة (arīka)}]{\ar{أريكة} (arīka) — a sofa}

\label{lesson:AR-C213-arika}



\begin{quote}
Before the new one: say the Arabic for transparent, then the Arabic for shiny.
\end{quote}

\subsection*{You'll want to know: \ar{أريكة}}
\textbf{\ar{أريكة}} — \emph{arīka} — \textquotedblleft{}a sofa\textquotedblright{}.

\begin{cousinweb}
\textbf{\ar{أريكة}} (\emph{arīka}) is built on the native Arabic root \emph{ʾ-r-k}, \textquotedblleft{}a couch\textquotedblright{} (Lane). Lane's \emph{arīka} is a couch within a curtained canopy.
\end{cousinweb}

\begin{grammarlens}[title={What you've built}]
One more word for this chapter.
\end{grammarlens}

\subsection*{Your turn}
\begin{itemize}
\raggedright
  \item \emph{Say it:} \emph{arīka}
  \item \emph{Say it:} \emph{arīka}, once more
  \item \emph{Say it:} say \emph{arīka}
  \item \emph{Recall:} say the Arabic for transparent, then the Arabic for shiny, then say \emph{arīka} again
\end{itemize}

\begin{tcolorbox}[breakable,colback=teal!4,colframe=teal!35!black,title={Before you move on}]
What does \ar{أريكة} mean? (\textquotedblleft{}A sofa\textquotedblright{}.) Say it once more.
\end{tcolorbox}
\section[\texorpdfstring{\ar{فرشاة} (furshā)}{فرشاة (furshā)}]{\ar{فرشاة} (furshā) — a brush}

\label{lesson:AR-C213-fursha}



\begin{quote}
Before the new one: say the Arabic for shiny, then the Arabic for a sofa.
\end{quote}

\subsection*{You'll want to know: \ar{فرشاة}}
\textbf{\ar{فرشاة}} — \emph{furshā} — \textquotedblleft{}a brush\textquotedblright{}.

\begin{cousinweb}
\textbf{\ar{فرشاة}} (\emph{furshā}) is built on the native Arabic root \emph{f-r-sh}, \textquotedblleft{}to spread out\textquotedblright{} (Lane). \emph{Furshāh} is a modern word, \textquotedblleft{}a thing for spreading\textquotedblright{}: a brush.
\end{cousinweb}

\begin{grammarlens}[title={What you've built}]
One more word for this chapter.
\end{grammarlens}

\subsection*{Your turn}
\begin{itemize}
\raggedright
  \item \emph{Say it:} \emph{furshā}
  \item \emph{Say it:} \emph{furshā}, once more
  \item \emph{Say it:} say \emph{furshā}
  \item \emph{Recall:} say the Arabic for shiny, then the Arabic for a sofa, then say \emph{furshā} again
\end{itemize}

\begin{tcolorbox}[breakable,colback=teal!4,colframe=teal!35!black,title={Before you move on}]
What does \ar{فرشاة} mean? (\textquotedblleft{}A brush\textquotedblright{}.) Say it once more.
\end{tcolorbox}
\section[\texorpdfstring{\ar{غطاء} (ghiṭāʾ)}{غطاء (ghiṭāʾ)}]{\ar{غطاء} (ghiṭāʾ) — a lid}

\label{lesson:AR-C213-ghita}



\begin{quote}
Before the new one: say the Arabic for a sofa, then the Arabic for a brush.
\end{quote}

\subsection*{You'll want to know: \ar{غطاء}}
\textbf{\ar{غطاء}} — \emph{ghiṭāʾ} — \textquotedblleft{}a lid\textquotedblright{}.

\begin{cousinweb}
\textbf{\ar{غطاء}} (\emph{ghiṭāʾ}) is built on the native Arabic root \emph{gh-ṭ-w}, \textquotedblleft{}to cover\textquotedblright{} (Lane). \emph{Ghiṭāʾ} is a covering.
\end{cousinweb}

\begin{grammarlens}[title={What you've built}]
One more word for this chapter.
\end{grammarlens}

\subsection*{Your turn}
\begin{itemize}
\raggedright
  \item \emph{Say it:} \emph{ghiṭāʾ}
  \item \emph{Say it:} \emph{ghiṭāʾ}, once more
  \item \emph{Say it:} say \emph{ghiṭāʾ}
  \item \emph{Recall:} say the Arabic for a sofa, then the Arabic for a brush, then say \emph{ghiṭāʾ} again
\end{itemize}

\begin{tcolorbox}[breakable,colback=teal!4,colframe=teal!35!black,title={Before you move on}]
What does \ar{غطاء} mean? (\textquotedblleft{}A lid\textquotedblright{}.) Say it once more.
\end{tcolorbox}
\section[\texorpdfstring{\ar{أثاث} (athāth)}{أثاث (athāth)}]{\ar{أثاث} (athāth) — furniture}

\label{lesson:AR-C213-athath}



\begin{quote}
Before the new one: say the Arabic for a brush, then the Arabic for a lid.
\end{quote}

\subsection*{You'll want to know: \ar{أثاث}}
\textbf{\ar{أثاث}} — \emph{athāth} — \textquotedblleft{}furniture\textquotedblright{}.

\begin{cousinweb}
\textbf{\ar{أثاث}} (\emph{athāth}) is built on the native Arabic root \emph{ʾ-th-th}, \textquotedblleft{}goods, furniture\textquotedblright{} (Lane). Lane's \emph{athāth} is household goods.
\end{cousinweb}

\begin{grammarlens}[title={What you've built}]
One more word for this chapter.
\end{grammarlens}

\subsection*{Your turn}
\begin{itemize}
\raggedright
  \item \emph{Say it:} \emph{athāth}
  \item \emph{Say it:} \emph{athāth}, once more
  \item \emph{Say it:} say \emph{athāth}
  \item \emph{Recall:} say the Arabic for a brush, then the Arabic for a lid, then say \emph{athāth} again
\end{itemize}

\begin{tcolorbox}[breakable,colback=teal!4,colframe=teal!35!black,title={Before you move on}]
What does \ar{أثاث} mean? (\textquotedblleft{}Furniture\textquotedblright{}.) Say it once more.
\end{tcolorbox}
\section[\texorpdfstring{\ar{غسيل} (ghasīl)}{غسيل (ghasīl)}]{\ar{غسيل} (ghasīl) — laundry, washing}

\label{lesson:AR-C213-ghasil}



\begin{quote}
Before the new one: say the Arabic for a lid, then the Arabic for furniture.
\end{quote}

\subsection*{You'll want to know: \ar{غسيل}}
\textbf{\ar{غسيل}} — \emph{ghasīl} — \textquotedblleft{}laundry, washing\textquotedblright{}.

\begin{cousinweb}
\textbf{\ar{غسيل}} (\emph{ghasīl}) is built on the native Arabic root \emph{gh-s-l}, \textquotedblleft{}to wash\textquotedblright{} (Lane). \emph{Ghasīl} is \textquotedblleft{}washed things\textquotedblright{}, the laundry.
\end{cousinweb}

\begin{grammarlens}[title={What you've built}]
One more word for this chapter.
\end{grammarlens}

\subsection*{Your turn}
\begin{itemize}
\raggedright
  \item \emph{Say it:} \emph{ghasīl}
  \item \emph{Say it:} \emph{ghasīl}, once more
  \item \emph{Say it:} say \emph{ghasīl}
  \item \emph{Recall:} say the Arabic for a lid, then the Arabic for furniture, then say \emph{ghasīl} again
\end{itemize}

\begin{tcolorbox}[breakable,colback=teal!4,colframe=teal!35!black,title={Before you move on}]
What does \ar{غسيل} mean? (\textquotedblleft{}Laundry, washing\textquotedblright{}.) Say it once more.
\end{tcolorbox}
